\begin{abstract}
Agents that revisit evolving data can repeat the same procedural mistake even as the underlying values change: using evidence released after a cutoff, confusing publication time with reporting time, or failing to ground a numerical change in contemporaneous documents. Reusable skills are a natural repair, but a standard targeted-versus-generic comparison can be misleading when the generic helper itself degrades the original agent. We ask when a reusable temporal procedure becomes an \emph{intervention-defined causal bottleneck}. On each frozen endpoint, we compare a targeted procedure (T), a complexity-matched target-blind helper (G), and the original no-skill agent (N); a procedure counts as a binding repair only when T beats both G and N. This criterion changes an observed scientific verdict: on a DeepSeek cutoff cell, T=100\% and G=72\%, but N=100\%, so the apparent +28-point two-arm gain is not a repair. We evaluate the criterion on 35 frozen endpoints reconstructed from first-party BEA, NOAA, and EIA release histories, totaling 1,326 retained runtime-valid rows. In non-ceiling regimes, documentary grounding reaches 60\% versus 20\% for both controls; held-out EIA cutoff filtering reaches 100\% versus 70\% (G) and 47.5\% (N), with endpoint sign test $p=0.0156$ for T--N; and held-out Kimi release alignment reaches 12/12 versus 4/12 (G) and 7/12 (N). Byte-identical procedures retain direction on compatible future releases, while cross-domain grounding is null and a second model is already at no-skill ceiling on the EIA cutoff tasks. The evidence therefore supports a conditional diagnostic of procedure bottlenecks, not universal skill gains or superiority of a skill container over an equivalent retrieval-side temporal implementation.
\end{abstract}
